\newif\ifglobal

%%%%%%%%%%%%%%%%%%%%%%%
%%% Compile to view %%%
%%%%%%%%%%%%%%%%%%%%%%%

%%%%%%%%%%%%%%%%%%%%%%%%%%%%%%%%%%%%%%%%%%%%%%%%%%%
%%% Readme:                                     %%%
%%% https://www.overleaf.com/read/bwdjvfjksvjy/ %%%
%%%%%%%%%%%%%%%%%%%%%%%%%%%%%%%%%%%%%%%%%%%%%%%%%%%

% >>> M?: marks a module setting number or important notice

% >>> M4: Set {./relative-path} to external files for this module <<<
\def \modulefiles {./12-BOOTCTRL}
% To add an external file, use:
% (Replace '---' with actual filename)
% '\input{\modulefiles/tables/---.tex}' for tables
% '\includegraphics{\modulefiles/figures/---}' for figures
% '\subfile{\modulefiles/---}' for register files

% >>> M1: This module is in global project? <<<
% 'YES' - uncomment; 'NO' - comment out
\globaltrue

\ifglobal
    \documentclass[main\_\_EN.tex]{subfiles}
   
\else
    \documentclass[openany, 10pt]{book}

    % >>> M2: In ./00-shared/preamble-trm-module, also find
    %         settings for visibility of todo notes, watermark <<<
    \usepackage{preamble-trm-module}

    % >>> M3: Temporary labels in English,
    %         you can add your own labels there <<<
    \myexternaldocument{./00-shared/temp-labels-en}

\fi %\ifglobal


\setcounter{secnumdepth}{4}


% >>> M5: If this module needs any updates in
% './00-shared/preamble-trm-module.sty'
% (include additional package, etc.), contact your module PM
% or create the file './00-shared/changelog.tex' make the
% necessary updates yourself and document them in this file <<<

% Make text left-aligned and allow latex to break pages naturally, comment out for full justification
\raggedright
\raggedbottom


\begin{document}


% Sets language into which pre-defined bits of text are translated
\selectlanguage{english}
\setenglishformatting

% Do not delete this block
\ifglobal
% Add the following front matter if
% this module is outside of global project
\else
    \InputIfFileExists{readme}{}{}
    \clearpage
    \listoftodos
    \listofdonetodos
    \tableofcontents
    %\listoftables
    %\listoffigures
\fi


%%%%%%%%%%%%%%%%%%%%%%%%%
%%% TRM Module Begins %%%
%%%%%%%%%%%%%%%%%%%%%%%%%

\hypertarget{bootctrl}{}
\chapter{Chip Boot Control}
\label{mod:bootctrl}


\section{Overview}
\label{sec:bootctrl-overview}

The chip allows for configuring the following boot parameters through strapping pins and eFuse parameters at power-up or a hardware reset, without microcontroller interaction.

\begin{itemize}
    \item \textbf{Chip boot mode}
    \begin{itemize}
        \item Strapping pins: GPIO60 and GPIO61
    \end{itemize}
    
    \item \textbf{VDD\_SPI voltage}
    \begin{itemize}
        \item Strapping pins: GPIO36
        \item eFuse parameters: EFUSE\_PMU\_FLASH\_POWER\_SEL and EFUSE\_PMU\_FLASH\_POWER\_SEL\_EN
    \end{itemize}

    \item \textbf{ROM messages printing}
    \begin{itemize}
        \item Strapping pin: GPIO60
        \item eFuse parameters: \hyperref[fielddesc:EFUSEUARTPRINTCONTROL]{EFUSE\_UART\_PRINT\_CONTROL} and \hyperref[fielddesc:EFUSEDISUSBSERIALJTAGROMPRINT]{EFUSE\_DIS\_USB\_SERIAL\_JTAG\_ROM\_PRINT}
    \end{itemize}
    
    \item \textbf{JTAG signal source}
    \begin{itemize}
        \item Strapping pin: GPIO37
        \item eFuse parameters: \hyperref[fielddesc:EFUSEDISPADJTAG]{EFUSE\_DIS\_PAD\_JTAG}, \hyperref[fielddesc:EFUSEDISUSBJTAG]{EFUSE\_DIS\_USB\_JTAG}, and \hyperref[fielddesc:EFUSEJTAGSELENABLE]{EFUSE\_JTAG\_SEL\_ENABLE}
    \end{itemize}
\end{itemize}

The default values of all the above eFuse parameters are 0, which means that they are not burnt. Given that eFuse is one-time programmable, once an eFuse bit is programmed to 1, it can never be reverted to 0. For how to program eFuse bits, please refer to Chapter \ref{mod:efuse} \textit{\nameref{mod:efuse}}.

During Chip Reset (see Chapter \ref{mod:resclk} \textit{\nameref{mod:resclk}}), hardware samples and stores the voltage level of strapping pins as a strapping bit of ``0" or ``1" in latches, and holds these bits until the chip is powered down or next chip reset. Software can read the latch status (strapping value) from \hyperref[fielddesc:GPIOSTRAPPING]{GPIO\_STRAPPING}.

\section{Functional Description}
\label{sec:bootctrl-func-descr}

This section introduces chip reset functions and the patterns of the strapping pins and eFuse values to invoke each function.

\begin{importantlisting}
Only documented patterns should be used. If an undocumented pattern is used, it may trigger unexpected behaviors.
\end{importantlisting}

\subsection{Default Configuration}

The default values of the strapping pins, namely the logic levels, are determined by pins' internal weak pull-up/pull-down resistors at reset if the pins are not connected to any circuit, or connected to an external high-impedance circuit.

\begin{table}[H]
    \centering
    \caption{Default Configuration of Strapping Pins}
    \label{tab:bootctrl-default-config-strapping-pins}
    \begin{tabular}{ | l | c | c | }
    \hline
        \rowcolor{lightgray}
        \textbf{Strapping Pin} & \textbf{Default Configuration} & \textbf{Bit Value} \\\hline
        \textbf{GPIO36}  & Floating & -- \\\hline
        \textbf{GPIO37}  & Floating & -- \\\hline
        \textbf{GPIO60}  & Weak pull-up & 1 \\\hline
        \textbf{GPIO61}  & Weak pull-up & 1 \\\hline
    \end{tabular}
\end{table}

To change the strapping bit values, apply external pull-down/pull-up resistors, or use host MCU GPIOs to control the voltage level of these pins when powering up \chipname{}. After the reset is released, the strapping pins work as normal-function pins.

\subsection{Chip Boot Mode Control}
\label{sec:bootModeControl}

GPIO60 and GPIO61 control the boot mode after the reset is released. See Table \ref{tab:bootctrl-boot-mode-control} \textit{\nameref{tab:bootctrl-boot-mode-control}}.

\begin{table}[H]
    \centering
    \caption{Boot Mode Control}
    \label{tab:bootctrl-boot-mode-control}

    \begin{threeparttable}
    \begin{tabular}{ | l | c | c | }
    \hline
        \rowcolor{lightgray}
        \textbf{Boot Mode} & \textbf{GPIO61} & \textbf{GPIO60} \\\hline
        
        % table contents
        \textbf{SPI Boot mode} \tnote{1} & \textbf{1} & Any value \\\hline
        Joint Download Boot mode \tnote{2} & 0 & 1 \\\hline
        Invalid combination\tnote{3} & 0 & 0 \\\hline 
        \end{tabular}

        \begin{tablenotes}
            \item[1] \textbf{Bold} marks the default value and configuration.
            \item[2] Joint Download Boot mode supports the following download methods:
                \begin{itemize}
                    \item USB-Serial-JTAG Download Boot
                    \item USB-OTG Download Boot
                    \item UART Download Boot
                    \item GPSPI Download Boot
                \end{itemize}
        \item[3]  Do not set both GPIO61 and GPIO60 to 0, or unexpected behavior will occur. 
        \end{tablenotes}
     \end{threeparttable}
\end{table}

In SPI Boot mode, the ROM bootloader loads and executes the program from SPI flash to boot the system. SPI Boot mode can be further classified as follows:

\begin{itemize}
    \item Normal Flash Boot: supports Secure Boot. The ROM bootloader loads the program from flash into SRAM and executes it. In most practical scenarios, this program is the 2nd stage bootloader, which later boots the target application.
    \item Direct Boot: does not support Secure Boot and programs run directly from flash. To enable this mode, make sure that the first two words of the bin file downloaded to flash are \texttt{0xaedb041d}. For more detailed process, see Figure \ref{fig:boot_flow} \textit{\nameref{fig:boot_flow}}.
\end{itemize}

In Joint Download Boot mode, users can download binary files into flash using UART0, USB Serial/JTAG, USB 2.0 High-Speed OTG, or the GPSPI slave interface. It is also possible to download binary files into SRAM and execute them from SRAM.

Figure \ref{fig:boot_flow} \textit{\nameref{fig:boot_flow}} shows the detailed boot flow of the chip.

\begin{figure}[H]
    \centering
    {\includegraphics[width=1.0\textwidth]{\modulefiles/figures/boot_flow.eps}}\\
    \caption{Chip Boot Flow}
    \label{fig:boot_flow}
\end{figure}

The following eFuse parameters control boot mode behaviors:

\begin{itemize}
    \item \hyperref[fielddesc:EFUSEDISFORCEDOWNLOAD]{\textbf{EFUSE\_DIS\_FORCE\_DOWNLOAD}}
    \begin{itemize}
    \item If this eFuse is 0 (default), software can force switch the chip from SPI Boot mode to Joint Download Boot mode by setting register \hyperref[fielddesc:LPAONFORCEDOWNLOADBOOT]{LP\_AON\_FORCE\_DOWNLOAD\_BOOT} and triggering a CPU reset. In this case, hardware overwrites \hyperref[fielddesc:GPIOSTRAPPING]{GPIO\_STRAPPING}[4:3] from ``1x" to ``01".

    \item If this eFuse is 1, \hyperref[fielddesc:LPAONFORCEDOWNLOADBOOT]{LP\_AON\_FORCE\_DOWNLOAD\_BOOT} is disabled, and \hyperref[fielddesc:GPIOSTRAPPING]{GPIO\_STRAPPING} cannot be overwritten.
    \end{itemize}

    \item \hyperref[fielddesc:EFUSEDISDOWNLOADMODE]{\textbf{EFUSE\_DIS\_DOWNLOAD\_MODE}}

    If this eFuse is 1, Joint Download Boot mode is disabled, and \hyperref[fielddesc:GPIOSTRAPPING]{GPIO\_STRAPPING} will not be overwritten by \hyperref[fielddesc:LPAONFORCEDOWNLOADBOOT]{LP\_AON\_FORCE\_DOWNLOAD\_BOOT}.

    \item \hyperref[fielddesc:EFUSEENABLESECURITYDOWNLOAD]{\textbf{EFUSE\_ENABLE\_SECURITY\_DOWNLOAD}}

    If this eFuse is 1, Joint Download Boot mode only allows writing and erasing plaintext flash and does not support any SRAM or register operations. Ignore this eFuse if Joint Download Boot mode is disabled.
    \item \hyperref[fielddesc:EFUSEDISDIRECTBOOT]{\textbf{EFUSE\_DIS\_DIRECT\_BOOT}}
    If this eFuse is 1, Direct Boot mode is disabled.
\end{itemize}

USB Serial/JTAG Controller can also force switch the chip to Joint Download Boot mode from SPI Boot mode, and vice versa. For detailed information, please refer to Chapter \ref{mod:usbserialjtag} \textit{\nameref{mod:usbserialjtag}}.

\subsubsection{Secure Debugging in Joint Download Boot Mode}
\label{sec:bootctrl-secure-debugging}

In Joint Download Boot mode, the chip supports commands related to the Secure Debug Controller (SDC) when eFuse field EFUSE\_SDC\_ENA is \texttt{2'b01} or \texttt{2'b10}. SDC can be started only in this mode.
Before use, burn the SDC public-key hash into eFuse. In Joint Download Boot mode, the ROM supports certificate request and certificate verification. The certificate is issued off chip, transferred to the chip, and verified in hardware.

After successful verification, SDC can latch the following outputs according to its configuration:

\begin{itemize}
    \item Re-enable JTAG. This output takes effect immediately.
    \item Re-enable download mode, even if \hyperref[fielddesc:EFUSEDISDOWNLOADMODE]{EFUSE\_DIS\_DOWNLOAD\_MODE} has been burned. This output takes effect at the next software reset.
    \item Force SPI Boot mode, so the application booted this way can be debugged. This output takes effect at the next software reset.
\end{itemize}

If not cleared earlier, these outputs remain until power-on reset. After power-on reset, repeat the SDC unlock flow to debug again.

For SDC commands, certificate format, and procedures, see Chapter \ref{mod:sdc} \textit{\nameref{mod:sdc}}.


\subsection{VDD\_SPI Voltage Control}
\label{sec:vdd-spi-voltage}

VDD\_SPI can be sourced from either the internal or external power supply:

\begin{itemize}
    \item \textbf{(Default) If register PMU\_ACTIVE\_VDD\_SPI\_PD\_EN is low, power is supplied by the internal POWER\_SWITCH}
    \item If register PMU\_ACTIVE\_VDD\_SPI\_PD\_EN is high, power is supplied externally
\end{itemize}


The actual VDD\_SPI supply voltage level is jointly determined by EFUSE\_PMU\_FLASH\_POWER\_SEL\_EN, GPIO36, and EFUSE\_PMU\_FLASH\_POWER\_SEL, as shown in Table \ref{tab:vdd-spi-voltage-control} \textit{\nameref{tab:vdd-spi-voltage-control}}.

\begin{table}[H]
    \centering
    \caption{VDD\_SPI Voltage Control}
    \label{tab:vdd-spi-voltage-control}

    \begin{tabular}{ | l | c | c | c |}
    \hline
    
    \rowcolor{lightgray}
    \textbf{Voltage Level}  & 
    \textbf{EFUSE\_PMU\_FLASH\_POWER\_SEL\_EN} & \textbf{GPIO36} & \textbf{EFUSE\_PMU\_FLASH\_POWER\_SEL} \\\hline
    
    \multirow{2}{*}{3.3 V} &
    0 & 1 & Ignored \\ \cline{2-4} &
    1 & Ignored &1  \\\hline
     \multirow{2}{*}{1.8 V} &
     0 & 0 & Ignored \\ \cline{2-4} &
     1 & Ignored & 0  \\\hline

    \end{tabular}
\end{table}

\subsection{ROM Messages Printing Control}

During the boot process, the messages by the ROM code can be printed to:

\begin{itemize}
    \item \textbf{(Default) UART0 and USB Serial/JTAG controller}
    \item USB Serial/JTAG controller
    \item UART0
\end{itemize}
    
\hyperref[fielddesc:EFUSEUARTPRINTCONTROL]{EFUSE\_UART\_PRINT\_CONTROL} and GPIO60 control ROM messages printing to \textbf{UART0} as shown in Table \ref{tab:uart-rom-msg-printing-control} \textit{\nameref{tab:uart-rom-msg-printing-control}}.

\begin{table}[H]
    \centering
    \caption{UART0 ROM Message Printing Control}
    \label{tab:uart-rom-msg-printing-control}

    \begin{threeparttable}

    \begin{tabular}{ | l | c | C{1.3cm} |}
    \hline
    
    \rowcolor{lightgray}
    \textbf{UART0 ROM Code Printing} & \textbf{\hyperref[fielddesc:EFUSEUARTPRINTCONTROL]{EFUSE\_UART\_PRINT\_CONTROL}} &   \textbf{GPIO60}  \\\hline

    \multirow{3}{*}{\textbf{Enabled}} & \textbf{0} & Ignored \\ \cline{2-3} 
     & 1 & 0 \\ \cline{2-3} 
     & 2 & 1 \\ \hline
    \multirow{3}{*}{Disabled} & 1 & 1 \\ \cline{2-3} 
     & 2 & 0 \\ \cline{2-3} 
     & 3 & Ignored \\ \hline
    
    \end{tabular}

        \begin{tablenotes}
            \item[1] \textbf{Bold} marks the default value and configuration.
        \end{tablenotes}
        
    \end{threeparttable}
\end{table}

\hyperref[fielddesc:EFUSEDISUSBSERIALJTAGROMPRINT]{EFUSE\_DIS\_USB\_SERIAL\_JTAG\_ROM\_PRINT} controls the printing to \textbf{USB Serial/JTAG controller} as shown in Table \ref{tab:usb-rom-msg-printing-control} \textit{\nameref{tab:usb-rom-msg-printing-control}}.

\begin{table}[H]
    \centering
    \caption{USB Serial/JTAG ROM Message Printing Control}
    \label{tab:usb-rom-msg-printing-control}

    \begin{threeparttable}

    \begin{tabular}{ | l | c |}
    \hline
    
    \rowcolor{lightgray}
    \textbf{USB Serial/JTAG ROM Code Printing} & \textbf{\hyperref[fielddesc:EFUSEDISUSBSERIALJTAGROMPRINT]{EFUSE\_DIS\_USB\_SERIAL\_JTAG\_ROM\_PRINT}} \\\hline

    \textbf{Enabled} & \textbf{0} \\\hline
    Disabled & 1 \\\hline
    
    \end{tabular}

        \begin{tablenotes}
            \item[1] \textbf{Bold} marks the default value and configuration.
        \end{tablenotes}
        
    \end{threeparttable}
\end{table}

\subsection{JTAG Signal Source Control}


The strapping pin GPIO37 can be used to control the source of JTAG signals during the early boot process. Before the strapping value is latched, the internal pull-up/pull-down resistors of this pin are not enabled, and the strapping value must be controlled by the external circuit that cannot be in a high-impedance state.

As Table \ref{tab:jtag-signal-src-control} \textit{\nameref{tab:jtag-signal-src-control}} shows, GPIO37 is used in combination with \hyperref[fielddesc:EFUSEDISPADJTAG]{EFUSE\_DIS\_PAD\_JTAG}, \hyperref[fielddesc:EFUSEDISUSBJTAG]{EFUSE\_DIS\_USB\_JTAG}, and \hyperref[fielddesc:EFUSEJTAGSELENABLE]{EFUSE\_JTAG\_SEL\_ENABLE}.


\setlength\tabcolsep{2pt}

\begin{table}[H]
    \centering
    \caption{JTAG Signal Source Control}
    \label{tab:jtag-signal-src-control}
    \begin{threeparttable}
    \begin{tabular}{ | l | c | c | c | c | }
    \hline
    
    \rowcolor{lightgray}
    \textbf{JTAG Signal Source} & \textbf{\hyperref[fielddesc:EFUSEDISPADJTAG]{EFUSE\_DIS\_PAD\_JTAG}} & \textbf{\hyperref[fielddesc:EFUSEDISUSBJTAG]{EFUSE\_DIS\_USB\_JTAG}} & \textbf{\hyperref[fielddesc:EFUSEJTAGSELENABLE]{EFUSE\_JTAG\_SEL\_ENABLE}} &   \textbf{GPIO37}       \\\hline

    \multirow{3}{*}{\textbf{USB Serial/JTAG Controller}} & \textbf{0} & \textbf{0} & \textbf{0} & Ignored \\ \cline{2-5} 
     & 0 & 0 & 1 & 1 \\ \cline{2-5} 
     & 1 & 0 & Ignored & Ignored \\ \hline
    \multirow{2}{*}{JTAG pins \tnote{2}} & 0 & 0 & 1 & 0 \\ \cline{2-5} 
     & 0 & 1 & Ignored & Ignored \\ \hline
    JTAG is disabled & 1 & 1 & Ignored & Ignored \\ \hline
    \end{tabular}

        \begin{tablenotes}
            \item[1] \textbf{Bold} marks the default value and configuration.
            \item[2] JTAG pins refer to MTDI, MTCK, MTMS, and MTDO.
        \end{tablenotes}
    \end{threeparttable}
\end{table}


\end{document}
